\cotefile{60}
\scanpage{60}
\mspage{4 - 34 -}
\hfill n° 392

entier sur $A$, il n'y a pas d'autre idéal que $0$ au-dessus de l'idéal $0$ de $A$, (chap. V, § 2, n° 1, cor. 1 de la prop. 1), donc il revient au même de dire que $\underline{a} \ne 0$ ou que $\underline{b} \ne 0$.

(i)\quad En vertu de ce qui précède, on peut se borner au cas où $R$ est un $B$-module de torsion. Si $\underline{b}$ est contenu dans un idéal premier $\underline{P} \in P(B)$, $\underline{a}$ est contenu dans $\underline{P} \cap A = \underline{p}$, qui est de hauteur $1$. Inversement, si $\underline{a}$ est contenu dans un idéal premier $\underline{p} \in P(A)$, il existe un idéal premier $\underline{P}$ de $B$ qui contient $\underline{b}$ et est au-dessus de $\underline{p}$, (chap. V, § 2, n° 1, cor. 2 du th. 1). L'assertion (1) résulte de ces remarques et du n° 4, prop. 9.

(ii)\quad Pour tout $B$-module de torsion de type fini $R$, posons $\varphi(R) = \chi_A(R_{[A]})$ ; il est clair que, (pour les $B$-modules de torsion de type fini), $\varphi$ vérifie les conditions 1) et 2) de la prop. 11 du n° 5, (compte tenu de (i)). Il existe par suite un homomorphisme $\vartheta : D(B) \longrightarrow D(A)$ tel que $\varphi(R) = \vartheta(\chi_B(R))$ pour tout $B$-module de torsion $R$ de type fini. L'homomorphisme $\vartheta$ est déterminé par sa valeur pour tout $B$-module de la forme $B/\underline{P}$ où $\underline{P} \in P(B)$, puisque $\chi_B(B/\underline{P}) = \underline{P}$. Or, pour tout idéal premier $\underline{q} \ne \underline{p} = \underline{P} \cap A$ dans $P(A)$, on a $\underline{q} \not\subset \underline{P}$, donc $(B/\underline{P})_{\underline{q}} = 0$. D'autre part, si on pose $S = A - \underline{p}$, $\underline{P}.S^{-1}B$ est un idéal maximal de $S^{-1}B$ et $(B/\underline{P})_{\underline{p}} = S^{-1}B/\underline{P}.S^{-1}B$ est isomorphe au corps des fractions de $B/\underline{P}$, c'est-à-dire au corps résiduel de $v_{\underline{P}}$ ; sa longueur en tant que $A_{\underline{p}}$-module est donc $f_{\underline{P}/\underline{p}}$ ; ce qui prouve que $\vartheta = N$, (n° 5, déf. 4).

(iii) Si $T$ est le sous-module de torsion de $R$, $T_{[A]}$ est le sous-module de torsion de $R_{[A]}$, et $(R/T)_{[A]} = R_{[A]}/T_{[A]}$ ; pour prouver (21), on peut donc se borner au cas où $R$ est sans torsion. Alors $R$ est identifié à un sous-$B$-module de $R_{(L)}$, et contient une base $(e_i)_{1\leq i\leq m}$ de $R_{(L)}$ sur $L$. Si $(b_j)_{1\leq j\leq n}$ est une base de $L$ sur $K$ formée d'éléments de $B$, les $b_j e_i$ constituent une base de $R_{(L)}$ sur $K$, formée d'éléments de $R$, d'où (21). Soit d'autre part $M$ un sous-$B$-module
% archive sheet number: 59

\scanpage{61}
Cohomologie \unsure{à} l'infini

$R^{\deleted{\unsure{\infty|*|i}}}_{\infty}f(F)$ pour $f$ \unsure{compactifiable} en $\bar f$

On \unsure{regarde|envisage} la suite exacte
\diagram{\begin{tikzcd}[ampersand replacement=\&,column sep=large]
\mathbb{R}i_{\deleted{\unsure{*}}!}(F) \arrow[r] \& \mathbb{R}\deleted{\unsure{f}}i_*(F) \arrow[r] \& \mathbb{R}i_{\unsure{*|\infty}}(F)\mid Y \arrow[ll,bend right=20]
\end{tikzcd}}
d'où en \unsure{appliquant} $\mathbb{R}\bar f_*$ une suite exacte
\diagram{\fbox{\begin{tikzcd}[ampersand replacement=\&,column sep=tiny,cramped]
\text{....} \arrow[r] \& \mathbb{R}^i f_!(F) \arrow[r] \& \mathbb{R}^i f_*(F) \arrow[r] \& R^i_{\infty}f(F) \arrow[r] \& \mathbb{R}^{i\unsure{-|+}1}f_!(F) \arrow[r] \& \unsure{\text{..}|\cdots}
\end{tikzcd}}}
en posant
\[
\mathbb{R}^i_{\infty}f(F)=\mathbb{R}^i\bar f_*(\mathbb{R}i_*(F)\mid Y)
\]

On peut aussi \unsure{considérer} la suite exacte \inserted{en \unsure{coh. relative}} pour $\deleted{\mathbb{R}i_*(F)}\inserted{\mathbb{R}i_!(F)}$ sur $\bar X$, et l'ouvert $\deleted{\unsure{Y}}X$ et $Y=\bar X-X$, \deleted{\unsure{on}} \deleted{\unsure{écrit}} \unsure{associée} \unsure{:}
\diagram{\begin{tikzcd}[ampersand replacement=\&,column sep=large]
\mathbb{R}j^!\unsure{(}\mathbb{R}i_!(F)) \arrow[r] \& \mathbb{R}i_!(F) \arrow[r,phantom,"\unsure{\text{---}|\longrightarrow}" description] \& \mathbb{R}i_*(F) \arrow[ll,bend right=20]
\end{tikzcd}}
\diagram{\begin{tikzcd}[ampersand replacement=\&,column sep=small,cramped]
\deleted{\mathbb{R}^i\unsure{-|\leftarrow}} \& \deleted{\mathbb{R}^i f_!(F)} \arrow[r,phantom,"\deleted{\longrightarrow}" description] \& \deleted{\mathbb{R}^i f_*(F)} \arrow[r,phantom,"\deleted{\longrightarrow}" description] \& {}
\end{tikzcd}}
et au décalage d'une unité près, $\mathbb{R}i_*(F)\mid Y$ est \unsure{figuré} par $\mathbb{R}j^!\mathbb{R}i_!(F)$ ---
% archive sheet number: 60

\scanpage{62}
\mspage{6 - 45 -}
\hfill n° 392

d) Déduire de c) un exemple d'anneau de valuation discrète $A$ et d'une extension séparable $L$ de $K$ tel que $B$ ne soit de la forme $A[x]$ pour aucun $x\in B$, (partir d'une extension cyclotomique de $\underline{Q}$, en utilisant l'exerc. 11 du chap. VI, § 8).

8)\qquad a) Soient $A$ un anneau de Dedekind, $K$ son corps des fractions, $L$ une extension finie séparable de $K$, $B$ la fermeture intégrale de $A$ dans $L$. On suppose que pour tout idéal premier $\underline{P}$ de $B$, le corps résiduel de $B$ en $\underline{P}$ soit séparable sur celui de $A$ en $\underline{P}\cap A$. Montrer qu'il existe un élément $x\in B$ tel que le $B$-module $\Omega_A(B)$ des différentielles de $B$ sur $A$ soit engendré par $dx$.

b) Donner un exemple d'un anneau de valuation discrète complet $A$, et d'une extension séparable $L$ du corps des fractions $K$ de $L$, tels que si $B$ est la fermeture intégrale de $A$ dans $L$, $\bar K$ et $\bar L$ les corps résiduels de $A$ et $B$ respectivement, $\bar L$ soit une extension radicielle non monogène de $\bar K$, (procéder comme dans le chap. VI, § 8, exerc. 10). Pour une telle extension, montrer que le $B$-module $\Omega_A(B)$ n'est pas monogène.

\begin{center}\texttt{-:-:-:-}\end{center}
% archive sheet number: 61

\scanpage{63}
$L(t)$ une \unsure{fonction} rat. en $t$, \deleted{\unsure{$q$ est}} $q=p^{\deleted{\unsure{f|a}}}$
\[
L\left(\frac{1}{qt}\right)=A(t)L(t) \qquad A(t)\text{ de ``niveau'' inférieur}
\]
On veut simplifier l'équation fonctionnelle, en \unsure{posant}
\[
L'(t)=\lambda(t)L(t)
\]
d'où
\[
\begin{aligned}
L'\left(\frac{1}{qt}\right)&=\lambda\left(\frac{1}{qt}\right)L\left(\frac{1}{qt}\right)=\lambda\left(\frac{1}{qt}\right)A(t)L(t)\\
&=\lambda\left(\frac{1}{qt}\right)\lambda(t)^{-1}A(t)\deleted{\unsure{\lambda}}L'(t)
\end{aligned}
\]
et on aimerait obtenir $L'\left(\deleted{\unsure{t}}\frac{1}{qt}\right)=c\deleted{\unsure{\lambda}}L'(t)$, i.e.
\[
\frac{\lambda\left(\frac{1}{qt}\right)}{\lambda(t)}=c\,A(t)^{-1}
\]

Si on trouve un tel $\lambda$, tout autre $\lambda'(t)=\alpha(t)\lambda(t)$ \inserted{avec une $\mathrm{c}^{\mathrm{te}}$ $c'$} est \unsure{déterminé} par \unsure{$\alpha$} \unsure{satisfaisant|solution générale}
\[
\frac{\alpha\left(\frac{1}{qt}\right)}{\alpha(t)}=\frac{c'}{c}\qquad \text{i.e.}\qquad \alpha\left(\frac{1}{qt}\right)=a\,\alpha(t)
\]
avec $a$ une $\mathrm{cte}$, \tikz[baseline=(v.base)]{\node[draw,ellipse,inner sep=1pt](v){néc égale à $\pm1$};}
\margin{donc dans l'équation \unsure{fonctionnelle} \unsure{modifiée} $c$ \unsure{est dét au signe}}
\par\noindent\rule{.72\linewidth}{.3pt}

Faisons opérer le groupe \deleted{\unsure{$G$|$\varphi$}} $\mathbb{Z}/2\mathbb{Z}$ sur les \inserted{\unsure{gpe mult.}} fonctions rationnelles \inserted{$\ne 0$}, via la substitution $t\longrightarrow\frac{1}{qt}$. Alors (pour $c$ donné) la question \unsure{en|de} \unsure{trouver} $\lambda$
% archive sheet number: 62

\scanpage{64}
\mspage{6 - 35 -}
\hfill n° 392

9) Soient $k$ un corps, $A$ l'anneau de polynômes $k[X,Y]$. Dans le' $A$-module $M=A^2$, on considère la forme linéaire $f$ telle que $f(e_1)=X$, $f(e_2)=Y$, ($(e_1,e_2)$ étant la base canonique). Montrer que le noyau $L$ de $f$ est un $A$-module libre monogène, mais que le quotient $M/L$, (qui est isomorphe à un idéal de $A$), n'est pas réflexif.

10) Soit $A$ un anneau commutatif. Pour tout sous-module $R$ d'un $A$-module libre de type fini $L=A^n$, on désigne par $\underline{c}_1(R)$ l'idéal engendré par les $\langle x,x^*\rangle$, où $x$ parcourt $R$ et $x^*$ parcourt le dual $L^*$ ; on pose $\underline{c}_k(R)=\underline{c}_1(\operatorname{Im}(\overset{k}{\bigwedge}R))$, $\operatorname{Im}(\overset{k}{\bigwedge}R)$ étant l'image canonique de la puissance extérieure $k$-ème de $R$ dans $\overset{k}{\bigwedge}L$, Si $R_1\subset R_2$ sont deux sous-modules de $L$, on a $\underline{c}_k(R_1)\subset \underline{c}_k(R_2)$ pour tout $k$.

a) Soit $M$ un $A$-module de type fini ; $M$ est isomorphe à un module quotient $L/R$, où $L=A^n$ pour un $n$ convenable. Montrer que la suite d'idéaux $(\underline{a}_k)_{k\geq0}$ telle que $\underline{a}_k\unsure{)|'}=\underline{c}_{n-k}\unsure{'}(R)$ pour $k\leq n$ et $\underline{a}_k=A$ pour $k>n$, est indépendante de l'expression de $M$ sous la forme $L/R$. (Considérer d'abord le cas où $L/R$ et $L/R'$ sont isomorphes ; remarquer ensuite que $A^n/R$ est isomorphe à $A^{n+h}/(R\times A^h)$ pour tout $h>0$). On pose $\underline{d}_k(M)=\underline{a}_k$ pour tout $k\geq0$, et on dit que ce sont les \uline{idéaux déterminantiels} associés à $M$. On a $\underline{d}_k(M)\subset \underline{d}_{k+1}(M)$ pour $k\geq0$.

b) Si $A$ est un anneau principal et $M=L/R$ où $L$ est libre de type fini, montrer que les idéaux $\underline{c}_{k+1}(R)(\underline{c}_k(R))^{-1}$ sont les \uline{facteurs invariants} de $R$ dans $L$.

c) Soit $r=\gamma(M)$ le plus petit des cardinaux des systèmes de générateurs de $M$, et soit $r_o$ le plus petit entier $h$ tel que $\underline{d}_h(M)=A$. Montrer que $r_o\leq r$ et donner un exemple où $r_o<r$, (prendre pour $A$ un anneau de Dedekind).

d) Si $\underline{a}$ est l'annulateur de $M$, montrer que l'on a $\underline{d}_o(M)\subset \underline{a}$ et $\underline{a}^{r-k}\subset \underline{d}_k(M)$ pour $k\leq r=\gamma(M)$.

e) On suppose que $M$ est somme directe de sous-modules $M_i$, $(1\leq i\leq h)$. Montrer que l'on a :
\[
\underline{d}_k(M)=\sum \underline{d}_{k_1}(M_1)\ldots \underline{d}_{k_h}(M_h)
\]
où la somme est étendue aux suites finies $(k_i)_{1\leq i\leq h}$ telles que $\sum_{i=1}^{h}k_i=k$.

f) Si $N$ est un sous-module de type fini de $M$, on a $\underline{d}_k(M/N)\subset \underline{d}_k(M)$ pour tout $k\geq0$. Montrer que l'on a :
\[
\sum_{j=o}^{k}\underline{d}_j(N)\underline{d}_{k-j}(M/N)\subset \underline{d}_k(M).
\]
% archive sheet number: 63

\scanpage{65}
\deleted{\unsure{équivaut} à celle \unsure{de} \unsure{nous} \unsure{on|en|un} \unsure{peut|faut|fct} \unsure{revenir}}
\[
\boxed{\underline{\lambda}^{g}\lambda^{-1}=}
\]

Il y a une solution triviale $\lambda=\frac{1}{L}$ i.e. $L'(t)=1$. Mais nous \unsure{désirons|devons} \deleted{\unsure{que}} \inserted{imposer} que $\lambda$ \deleted{\unsure{$L$}} soit de ``niveau'' \unsure{strictmt|strict.} inf à $L$.

Pour préciser ce point, nous \unsure{supposons},

\deleted{$L(0)=1$, que les racines et pôles \unsure{de} $L$ sont en valeurs absolues de la forme $p^{-\frac{i}{2}}$ avec $\deleted{2}\rho\leq i\leq\deleted{2}\rho+2n$ i.e. $\rho\leq\frac{i}{2}\leq\rho+n$, \deleted{et \unsure{modulo multiplication|multipliant|multipliées|modulo coefficients} par une \unsure{puissance} de $p^{1/2}$, \unsure{des|dans} des valeurs \unsure{propres|propre} \unsure{entières}} et on veut se restreindre aux $\lambda$ ayant \uline{la $=$ propriété}, \deleted{et de plus} (de poids compris entre $\rho$ et $\rho+n$) et de plus de ``niveau'' $\leq\deleted{\unsure{\rho+n+1}}$}
% margin note written sideways
\margin{et sont \unsure{tous|toutes} de la forme $(p^{1/2})^{\deleted{\unsure{i}}\unsure{\nu}}\zeta_i$, où $\nu$ un entier $\zeta_i$ un \uline{entier alg.}}
\[
L(t)=\underbrace{L_{\rho}(t)\cdots L_{\rho\deleted{\unsure{b}}+n}(t)}_{\unsure{\longrightarrow|1|\ne}}L'_{\rho+n-1}(pt)L'_{\rho+n-2}(p^2t)\cdots L'_{\rho}(p^nt)
\]
où les valeurs propres inverses des $L_{\deleted{\unsure{2}}i}$ sont \deleted{\unsure{et des}} des entiers alg en val. absolue $p^{i/2}$, et de $=$ pour $L'_i$ (donc $L'_{\rho+n-i}(p^it)$ \unsure{à|a} des valeurs propres inverses qui sont des entiers alg de val. absolue $p^{\frac{\rho+n-i}{2}}p^i=p^{\frac{\rho+n+i}{2}}$.
% archive sheet number: 64

\scanpage{66}
\mspage{6 - 34 -}
\hfill n° 392

6) Soient $k$ un corps, $B$ l'anneau de polynômes $k[X,Y]$, $A$ le sous-anneau $k[X^2,XY,Y^2]$ de $B$.

a) Montrer que $A$ est un anneau noethérien intégralement clos et que $B$ est un $A$-module de type fini.

b) Montrer que l'idéal $\underline{p}=AX^2+AXY$ de $A$ est un idéal premier de hauteur 1, (donc divisoriel), mais que le $B$-module $\underline{p}\otimes_A B$ a un module de torsion $\ne 0$ et n'est donc pas réflexif ; l'idéal $\underline{p} B$ de $B$ n'est pas divisoriel, l'application canonique $\underline{p}\otimes_A B\longrightarrow \underline{p} B$ n'est pas injective et $\underline{p}$ n'est pas un $A$-module plat.

7) Soient $A$ un anneau noethérien intégralement clos, $E$ un $A$-module sans torsion de type fini, $E^*$ son dual.

a) Montrer que l'homomorphisme canonique $E^*\otimes_A E\longrightarrow\operatorname{End}_A(E)$ est un pseudo-isomorphisme.

b) Déduire de a) que pour que $E$ soit un $A$-module projectif, il faut et il suffit que $E^*\otimes_A E$ soit un $A$-module réflexif. (Remarquer que si $E^*\otimes_A E$ est réflexif, l'homomorphisme canonique $E^*\otimes_A E\longrightarrow\operatorname{End}_A(E)$ est bijectif).

\textsuperscript{*}8) Soient $A$ un anneau \unsure{npethérien|noethérien} intégralement clos, $M_1$, $M_2$ deux $A$-modules de type fini.

a) Montrer que les $A$-modules $\operatorname{Tor}_i^A(M_1,M_2)$, $\operatorname{Ext}_A^i(M_1,M_2)$ sont pseudo-nuls pour $i\geq2$, (se ramener au cas où $A$ est principal).

b) Si $M_1$ est sans torsion, montrer que $\operatorname{Tor}_1^A(M_1,M_2)$ et $\operatorname{Ext}_A^1(M_1,M_2)$ sont pseudo-nuls, (même méthode).

c) Si $M_1$ est un $A$-module de torsion, montrer que l'on a, (avec les notations du $n^{\circ}$ \quad) :
\[
\begin{aligned}
\chi(M_1\otimes_A M_2)-\chi(\operatorname{Tor}_1^A(M_1,M_2))&=r(M_2)\,\chi(M_1)\\
\chi(\operatorname{Hom}_A(M_1,M_2))-\chi(\operatorname{Ext}_A^1(M_1,M_2))&=-r(M_2)\,\chi(M_1)\\
\chi(\operatorname{Hom}_A(M_2,M_1))-\chi(\operatorname{Ext}_A^1(M_2,M_1))&=r(M_2)\,\chi(M_1)
\end{aligned}
\]

d) Montrer que pour tout couple de $A$-modules de type fini $M_1$, $M_2$, on a :
\[
\begin{aligned}
c(M_1\otimes_A M_2)-c(\operatorname{Tor}_1^A(M_1,M_2))&=r(M_1)c(M_2)+r(M_2)c(M_1)\\
c(\operatorname{Hom}_A(M_1,M_2))-c(\operatorname{Ext}_A^1(M_1,M_2))&=r(M_1)c(M_2)-r(M_2)c(M_1).\;*
\end{aligned}
\]
% archive sheet number: 65

\scanpage{67}
Ceci posé, on \deleted{\unsure{veut se}} se permet de ``corriger'' par des $\lambda$ \unsure{seulmt|seuls|soumis} de la forme
\[
\begin{aligned}
&\lambda_{\rho}(t)\cdots\lambda_{\rho+n\deleted{\unsure{+}}-1}(t)\,\lambda_{\rho+n-2}(pt)\lambda_{\rho+n-1}(pt)\lambda_{\rho+n-2}(p^2t)\\
&\hspace{15em}\cdots\lambda_{\rho}(p^nt)
\end{aligned}
\]
donc de poids entre $\rho$ et $\rho+n$ comme $L$, mais de \uline{niveau} $\leq\rho+n-1$ (et non \unsure{seulmt|seulement} $\leq\rho+n$ comme $L$)

\deleted{Si le \unsure{facteur|terme} $L_{\rho+n}$}

Mais \unsure{supposons} de plus que $A(t)$ lui-$=$ est de \deleted{\unsure{niv}} ce type.
\[
\begin{aligned}
L\left(\frac{1}{p^{n+\rho}t}\right)
 &=L'_{\rho}\left(\frac{1}{p^\rho t}\right)L'_{\rho+1}\left(\frac{1}{p^{\rho+1}t}\right)\cdots L'_{\rho+n-1}\left(\frac{1}{p^{\unsure{n+\rho-1}}t}\right)\\
 &\qquad L_{\rho+n}\left(\frac{1}{p^{\rho+n}t}\right)L_{\rho+n-1}\left(\frac{1}{p^{\unsure{n+\rho}}t}\right)\cdots L_{\rho}\left(\frac{1}{p^{\unsure{n+\rho}}t}\right)
\end{aligned}
\]
De plus, on suppose que \unsure{chaque} $L_i$ et $L'_i$ a une équation fonctionnelle du type \unsure{correspondant}
\[
\begin{aligned}
&=\delta'_{\rho}t^{\chi'_{\rho}}L'_{\rho}(t)\cdots\delta'_{\unsure{\rho+n-1}}t^{\chi'_{\rho+n-1}}L'_{\rho+n-\unsure{i|1}}(t)\\
&\qquad \deleted{\unsure{L_t|L}}\delta_{\rho+n}t^{\chi_{\rho}}L_{\rho+n}(t)\,\delta_{\rho+n-1}t^{\chi_{\rho+n-1}}L_{\rho+n-1}(pt)\\
&\hspace{18em}\unsure{\delta|\Delta}_{\rho}t^{\chi_{\rho}}L_{\rho}(p^nt)\\
&=\delta t^\chi L'_{\rho}(t)\cdots L'_{\rho+n-1}(t)L_{\rho+n}(t)L_{\rho+n-1}(\deleted{\unsure{q}}pt)\cdots L_{\rho}(p^nt)
\end{aligned}
\]
% archive sheet number: 66

\scanpage{68}
\mspage{6 - 33 -}
\hfill n° 392

\S~4.

1) Soient $A$ un anneau noethérien intégralement clos, $V$ un espace vectoriel de dimension finie sur le corps des fractions de $A$. Montrer que si $(M_\lambda)$ est une famille quelconque de réseaux réflexifs de $V$ contenant toutes un même réseau $N$, le réseau $M=\bigcap_\lambda M_\lambda$ est réflexif, (considérer les réseaux associés $M_\lambda^*$).

\textsuperscript{*}2) Soient $A$ un anneau noethérien intégralement clos, $E$, $F$ deux $A$-modules de type fini. On suppose que $E$ est sans torsion, que $\operatorname{Hom}_A(E,F)$ est réflexif et que $\operatorname{Ext}_A^1(E,F)=0$. Montrer alors que $F$ est réflexif. (Prouver d'abord que $F$ est sans torsion ; si $T=F^{**}/c_F(F)$, calculer de deux manières $\operatorname{Ass}(\operatorname{Hom}_A(E,T))$ en utilisant la suite exacte :
\diagram{\begin{tikzcd}[ampersand replacement=\&,column sep=small]
0 \arrow[r] \& \operatorname{Hom}(E,F) \arrow[r] \& \operatorname{Hom}(E,F^{**}) \arrow[r] \& \operatorname{Hom}(E,T) \arrow[r] \& 0
\end{tikzcd}}
et le chap. IV, \S~1, $n^{\circ}$ \quad, prop. \quad). \textsubscript{*}

3) Soient $A$ un anneau noethérien intégralement clos, $K$ son corps des fractions, $M$ un réseau réflexif d'un espace vectoriel $V$ de dimension finie sur $K$, $L$ un réseau \uline{libre} de $V$ contenant $M$. Montrer qu'il existe un réseau libre $L_1$ de $V$ tel que $M=L\cap L_1$. (Considérer l'ensemble fini $I$ des idéaux premiers $\underline{p}$ de hauteur 1 de $A$ tels que $L_{\underline{p}}\ne M_{\underline{p}}$, et l'anneau principal $S^{-1}A$, où $S=\bigcap_{\underline{p}\in I}(A-\underline{p})$ ; montrer qu'il existe un réseau libre $L_o$ de $V$ tel que $M_{\underline{p}}=(L_o)_{\underline{p}}$ pour tout $\underline{p}\in I$, et un $s\in S$ tel que $M\subset s^{-1}L_o=L_1$).

4) Soient $k$ un corps, $A$ l'anneau de polynômes $k[X,Y]$.

a) Soit $(e_1,e_2)$ la base canonique du $A$-module $A^2$, et soit $E$ le sous-$A$-module de $A^2$ engendré par $(X-Y)e_1$, $e_1+Xe_2$ et $e_1+Ye_2$ ; soit $F$ le sous-module monogène de $E$ engendré par $(X-Y)^2e_1$. Montrer que le $A$-module $M=E/F$ n'est pas somme directe de son sous-module de torsion et d'un module sans torsion.

b) Montrer que le $A$-module de torsion $A/AXY$ n'est pas somme directe de sous-modules monogènes de la forme $A/\underline{P}_i^{n_i}$, où les $\underline{P}_i$ sont des idéaux premiers de hauteur 1.

5) Soient $A$ un anneau noethérien intégralement clos, $M$ un $A$-module de type fini. Montrer que si $\underline{a}$, $\underline{b}$ sont deux idéaux de $A$, le $A$-module $((\underline{a} M)\cap(\underline{b} M))/(\underline{a}\cap\underline{b})M$ est pseudo-nul ; donner un exemple où il n'est pas nul.
% archive sheet number: 67

\scanpage{69}
\deleted{\unsure{Au point|Au fait}}

\[
\frac{L\left(\frac{1}{qt}\right)}{L(t)}=\delta t^\chi
\left[\frac{L'_\rho(t)}{L'_\rho(p^nt)}\overset{\unsure{\text{--}}}{\cdots}
\frac{L'_{\rho+n-1}(t)}{L'_{\rho+n-1}(pt)}\right]
\left[\frac{L_{\rho+n-1}(pt)}{L_{\rho\deleted{\unsure{t}}+n-1}(t)}
\deleted{\unsure{L_\rho|L}}\cdots
\frac{L_\rho(p^nt)}{L_\rho(t)}\right]
\]
\[
A(t)=\delta t^\chi\prod_i\frac{L'_i(t)}{L'_i(p^{\rho+n-i}t)}
\prod_i\frac{L_{\deleted{\unsure{\rho}}i}(\deleted{\unsure{t|p}}p^{\rho+n-i}t)}{L_i^{\unsure{\prime|{}}}(t)}
\]

\[
\begin{aligned}
n=1\qquad &\underset{\lfloor\!\rule{2.8em}{.3pt}\!\rfloor}{L_\rho(t)}\,L_{\rho+1}(t)\,\underset{\lfloor\!\rule{3.2em}{.3pt}\!\rfloor}{L'_\rho(pt)}\\
&L_\rho(t)L'_\rho(pt)
\end{aligned}
\]
% archive sheet number: 68

\scanpage{70}
\mspage{6 - 32 -}
\hfill n° 392

26) Soient $A$ un anneau factoriel, $K$ son corps des fractions. Pour que des éléments $f_i$, $(1\leq i\leq r)$, de $B=A[X_1,\ldots,X_n]$ soient tels que l'idéal $\sum_{i=1}^{r}Bf_i$ soit égal à $B$, il faut et il suffit qu'il existe des polynômes $v_i$, $(1\leq i\leq r)$, de $K[X_1,\ldots,X_n]$ tels que $\sum_{i=1}^{r}v_i f_i=1$ et qui vérifient en outre la condition suivante : si l'on pose $v_i=w_i/d$, où $d\in A$, les polynômes $w_i$ appartiennent à $A[X_1,\ldots,X_n]$ et le p. g. c. d. de l'ensemble des coefficients de tous les $w_i$, $(1\leq i\leq r)$, est égal à 1, alors, pour tout élément extrémal $p$ de $A$ divisant $d$, l'idéal engendré par les classes des $f_i$ dans l'anneau $(A/Ap)[X_1,\ldots,X_n]$ est cet anneau tout entier.

\begin{center}
-:-:-:-
\end{center}
% archive sheet number: 69
